\section{Technical appendix}\label{app:technical}

The main text keeps the proof sketches close to the economic argument. This appendix
records the two mathematical details from the longer draft that are still useful for a
research-grade version of the paper: why the option-only optimizer has a well-defined
scale only after operational budgets are imposed, and how the maximum-Sharpe solver is
connected to a convex conic program.

\subsection{Scale, costs, and no-free-exposure logic}

\begin{proposition}[No free option exposure]\label{prop:app_no_free_exposure}
Let \(\mathcal C\) include binding gross-premium, short-premium, stress-loss,
margin/collateral, or contract-level liquidity controls. Then no economically nonzero
sequence of option positions can take unbounded premium, payoff, or Greek exposure while
remaining feasible and using zero capital.
\end{proposition}

\begin{proof}
The listed controls are positively homogeneous in the premium-weight vector. If
\(\|w_k\|\to\infty\) along an economically nonzero direction, at least one controlled
premium, margin, stress, Greek, or liquidity functional diverges. A direction in the null
space of all such functionals has no priced premium exposure and no admissible tradable
contract size in this model. Therefore every economically meaningful option direction
must consume a real budget before scale can diverge.
\end{proof}

This proposition is the formal reason the paper does not treat gross option Sharpe as a
trading claim. A cheap option can carry large return exposure per dollar of premium, and
a short option can receive premium while creating a future liability. The optimizer is
only an allocation rule after premium, short-premium, stress, margin, concentration, and
liquidity limits turn the scale-free return object into a feasible book.

\begin{proposition}[Operational monotonicity]\label{prop:app_monotonicity}
For
\[
  V(\mathcal C,c)=
  \max_{w\in\mathcal C}
  \widehat\mu^\top w-\frac{\gamma}{2}w^\top\Sigma w-c^\top |w|,
\]
with \(\Sigma\succeq0\), if \(\mathcal C_1\subseteq\mathcal C_0\) and
\(c_1\ge c_0\) componentwise, then \(V(\mathcal C_1,c_1)\le V(\mathcal C_0,c_0)\).
\end{proposition}

\begin{proof}
For every \(w\in\mathcal C_1\), the same \(w\) is feasible for \(\mathcal C_0\), and
\(-c_1^\top |w|\le -c_0^\top |w|\). Maximizing a weakly lower objective over a weakly
smaller feasible set cannot increase the value.
\end{proof}

The monotonicity result justifies the conservative reporting convention. Higher spread,
fee, slippage, margin, assignment, or impact assumptions can only make the utility value
worse. Tighter displayed-size, stress, or margin budgets can only make the feasible set
smaller. Positive results under full costs and hard caps are therefore more meaningful
than positive gross results, while failures under exact execution data would supersede the
current proxy-spread sensitivity.

\subsection{Mean estimation error, shrinkage, and purge details}\label{app:mean_error}

\begin{proof}[Proof of Proposition~\ref{prop:short_mean_error}]
Each error \(e_i\) is sub-Gaussian with variance proxy \(\sigma_i^2/T\le\bar\sigma^2/T\),
so for every \(\lambda>0\),
\(\E\exp(\lambda e_i)\le\exp(\lambda^2\bar\sigma^2/2T)\) and the same holds for \(-e_i\).
Let \(M=\max_i|e_i|\), the maximum of the \(m=2n\) variables \(\{\pm e_i\}\). By Jensen's
inequality and a union bound on the moment generating function,
\[
  \exp\big(\lambda\,\E[M]\big)\le\E\exp(\lambda M)
  \le\sum_{j=1}^{m}\E\exp(\lambda X_j)
  \le m\exp\!\Big(\frac{\lambda^2\bar\sigma^2}{2T}\Big),
\]
so \(\E[M]\le\log(m)/\lambda+\lambda\bar\sigma^2/2T\). Optimizing at
\(\lambda=\sqrt{2T\log m}/\bar\sigma\) gives
\(\E[M]\le\bar\sigma\sqrt{2\log(2n)/T}\), the first display. For the second, H\"older's
inequality gives \(\sup_{w\in\mathcal C}e^\top w\le\sup_{\|w\|_1\le g}e^\top w
= g\|e\|_\infty=gM\); take expectations. For the objective-loss statement, write the
optimizer's concave criterion as \(F_\mu(w)=\mu^\top w-R(w)\), where \(R\) collects the
risk and cost penalties, and let \(\widetilde w\) maximize \(F_{\widehat\mu}\) and
\(w^\ast\) maximize \(F_{\mu}\) over \(\mathcal C\). Then
\[
  F_\mu(\widetilde w)
  \;\ge\; F_{\widehat\mu}(\widetilde w)-g\|e\|_\infty
  \;\ge\; F_{\widehat\mu}(w^\ast)-g\|e\|_\infty
  \;\ge\; F_{\mu}(w^\ast)-2g\|e\|_\infty,
\]
so the expected sacrifice of true objective value is at most
\(2g\,\E\|e\|_\infty\le2g\bar\sigma\sqrt{2\log(2n)/T}\).
\end{proof}

\begin{proof}[Proof of Corollary~\ref{cor:short_crossover}]
The sandwich argument above holds verbatim when the mean input is the structural
forecast: with deterministic bias vector \(b=\mu^{\mathrm{struct}}-\mu\) and
\(\|b\|_\infty\le\beta\), the loss bound is \(2g\beta\). Comparing the two worst-case
bounds, the historical mean has the smaller guarantee exactly when
\(\bar\sigma\sqrt{2\log(2n)/T}<\beta\), which rearranges to
\(2\log(2n)<T\beta^2/\bar\sigma^2\), i.e.
\(n<\tfrac12\exp(T\beta^2/2\bar\sigma^2)\). At fixed \(T\) the left side grows in \(n\)
while the right side is constant, which yields the stated crossover.
\end{proof}

The bounds compare worst-case guarantees, not realized losses, and the sub-Gaussian model
ignores the correlation and missingness of real option panels; both simplifications favor
the historical mean, so the crossover argument is conservative in the direction the paper
uses it.

\begin{lemma}[One-month purge suffices for monthly option labels]\label{lem:purge_gap}
Index each label by its return month-end. Assume every tested contract's information
window spans at most \(\Delta\) calendar days and ends at its return month-end, where
\(\Delta\approx44\) days is the maximum decision-to-expiry span of the monthly-tenor
contracts. If the realized calendar gap between the latest training return month-end and
the earliest test return month-end is \(G>\Delta\) in every fold, then no training
label's information window overlaps any test label's window.
\end{lemma}

\begin{proof}
A training label with return month-end \(t_r\) has window contained in
\([t_r-\Delta,\,t_r]\). A test label with return month-end \(t_e\) has window contained
in \([t_e-\Delta,\,t_e]\), which begins at \(t_e-\Delta\). Overlap requires
\(t_r>t_e-\Delta\), i.e. \(t_e-t_r<\Delta\). The purge/embargo schedule enforces
\(t_e-t_r\ge G>\Delta\) for every training/test pair, so the windows are disjoint.
\end{proof}

The quantity \(G\) is not assumed; it is measured. The verifier recomputes the minimum
gap across all folds, schemes, and configurations from the fold schedule and ledger
artifacts and asserts \(G>44\) days; in the current artifacts the realized minimum is
\(G=59\) days. The one-month purge plus one-month embargo is therefore sufficient with
two weeks of slack, and the sufficiency claim is re-audited on every verification run
rather than hand-typed.

\subsection{Maximum Sharpe and conic reduction}

Let \(\mu=\widehat\mu_t\) and \(\Sigma=\widehat\Sigma_{O,t}\). When
\(\Sigma\succ0\), the unconstrained maximum-Sharpe direction is the usual tangency
direction even though the assets are option cashflows:
\[
  \argmax_{w\ne0}\frac{w^\top\mu}{\sqrt{w^\top\Sigma w}}
  \quad \propto \quad
  \Sigma^{-1}\mu .
\]
The difference from stock Markowitz is not the algebra of the tangency direction; it is
the construction of \(\mu\), \(\Sigma\), and the feasible set. Option premium weights have
cash and collateral as the numeraire, so the unit-stock-budget equation is replaced by
gross premium, short-premium, stress, margin, and liquidity budgets.

\begin{proposition}[Conic maximum-Sharpe form]\label{prop:app_conic}
Let \(K\subseteq\R^n\) be a closed convex cone containing some \(w\) with
\(w^\top\mu>0\), and let \(\Sigma\succ0\). Then
\[
  \sup_{w\in K\setminus\{0\}}
  \frac{w^\top\mu}{\sqrt{w^\top\Sigma w}}
  =
  \max\{\,y^\top\mu:\ y\in K,\ y^\top\Sigma y\le1\,\}.
\]
The right side is a convex quadratically constrained program and is a second-order-cone
program when a Cholesky or square-root factor of \(\Sigma\) is used.
\end{proposition}

\begin{proof}
For any nonzero \(w\in K\), define \(y=w/\sqrt{w^\top\Sigma w}\). Since \(K\) is a cone,
\(y\in K\), \(y^\top\Sigma y=1\), and \(y^\top\mu\) equals the Sharpe ratio of \(w\). This
shows that the right side is at least the left. Conversely, any feasible \(y\) with
\(y^\top\mu>0\) satisfies
\[
  \frac{y^\top\mu}{\sqrt{y^\top\Sigma y}}\ge y^\top\mu,
\]
so the left side is at least the right. Positive definiteness makes the risk ball
bounded, and the objective is linear on a closed convex feasible set.
\end{proof}

The production-candidate problem in the paper has bounded budgets rather than a pure
cone, so implementation uses the same geometry with a gross-risk or gross-premium
normalization and then imposes the operational caps directly. Absolute costs are handled
by positive and negative parts, \(w=w^+-w^-\), but the liquidity cap is imposed on the
net economic variable:
\[
  -\bar w_{i,t}\le w_i^+-w_i^-\le \bar w_{i,t}.
\]
Imposing the cap separately on \(w_i^+\) and \(w_i^-\) is not equivalent. It allows
offsetting positive legs to satisfy an algebraic representation while consuming displayed
liquidity without taking the intended net contract exposure.

\subsection{Risk-model estimator repair}

The population covariance identity in Proposition~\ref{prop:short_covariance} is
positive semidefinite by construction. The finite-sample estimate can still fail PSD in
one practical place: residual covariance estimated from pairwise-complete option
residuals on an unbalanced contract panel. The implementation therefore applies the
following repair before the matrix can enter an optimizer:
\[
  \widehat\Sigma_{\varepsilon,t}=Q\Lambda Q^\top
  \quad\Longrightarrow\quad
  \widehat\Sigma_{\varepsilon,t}^{+}
  =Q\max(\Lambda,\lambda_{\min}I)Q^\top ,
\]
with a small nonnegative eigenvalue floor. The factor term
\(\widehat B_t\widehat\Omega_t\widehat B_t^\top\) is PSD whenever
\(\widehat\Omega_t\succeq0\), so the repair is targeted to the residual block rather than
used as a blanket way to hide a bad risk model.

E1 then makes one additional research choice: it keeps the residual covariance diagonal.
That is not a statement that residual option risks are literally independent. It is a
regularization choice motivated by the breadth experiment. In the broad universe, dense
residual covariance estimated from short, changing contract panels became a false
diversification channel. Diagonal residual risk preserves contract-specific unexplained
variance while refusing to invert noisy off-diagonal residual moments.

\subsection{Appendix use and audit trail}

Mathematical mechanics that strengthen the results but would slow the main reading path
are kept in this appendix. Large empirical ledgers remain machine-readable artifacts
under \texttt{analysis/artifacts/breadth\_solutions/robustness/}; the body prints only
the scoreboard, distributional summaries, robustness heatmap, capacity/spread panel, and
rolling OOS paths. Every number in the body traces to one of those ledgers, and the
verification harness recomputes the headline statistics from the artifacts on each
build.
